% !TEX program = lualatexmk
% !TEX encoding = UTF-8 Unicode

\documentclass[12pt]{book}
\usepackage{graphicx} % for logo
\usepackage[language=english]{lipsum}   % for dummy text

\DeclareRobustCommand{\LaPhYs}
  {%
    \mbox{%
    L%
    %\kern-0.36em\raise0.8ex\hbox{\textsubscript{A}}%
    %\kern-0.36em\raise0.8ex\hbox{\textsubscript{A}}%
    \kern-0.36em\raisebox{0.8ex}{\textsubscript{A}}%
    \kern-0.15em P%
    %\kern-0.29em\lower1.0ex\hbox{H}%
    %\kern-0.29em\lower0.4ex\hbox{\scalebox{0.9}{H}}%
    \kern-0.29em\raisebox{-0.4ex}{\scalebox{0.9}{H}}%
    \kern-0.26em Y%
    %\kern-0.25em\raise2.0ex\rotatebox[origin=c]{-20}{\textsc{s}}%
    %\kern-0.25em\raisebox{\depth}{\rotatebox[origin=c]{-20}{\textsc{s}}}%
    %\kern-0.25em\raisebox{0ex}{\rotatebox[origin=c]{-28}{\textsc{s}}}%
    \kern-0.25em\rotatebox[origin=c]{-28}{\textsc{s}}%
    }%
  }%

\title{\LaPhYs: \LaTeX{} for Introductory Physics \\ (Working Draft)}
\author{Paul J. Heafner}

\begin{document}
\frontmatter

\maketitle
% !TEX program = lualatexmk
% !TEX encoding = UTF-8 Unicode

\chapter{Preface}
This book is for students taking introductory physics courses that require
traditional, written problem solutions.
\marginpar{\raggedright The \LaPhYs{} logo pokes fun at the spirit of this project.} 
It is also for instructors who teachor oversee such courses.\footnote{The \LaPhYs{} 
logo pokes fun at the spirit of this project.}
The \LaPhYs{} logo pokes fun at the spirit of this 
project. When students take humanities courses, they are almost always 
required to learn and use certain formats and citation schemes for 
papers to be turned in. STEM courses are different in that such papers
are rarely required, and when they are most instructor ask students to use a
traditional word processor like Microsoft Word. Some students will need to
learn to use Word and even students who have used it before may not be
familiar with the latest version or the version forced upon them by their
institutions. Even for those of us who witnessed the rise of traditional
word processors to ubiquity (who remembers WordStar or WordPerfect?) using 
them nowadays is difficult because menu organization often changes from one
version to another. Once students graduate or otherwise leave academia, 
student and academic discounts for commercial software typically expire and
students can no longer freely access those products.  In STEM course we
like to blah blah...

This book is also for physics instructors and faculty who may not be up to
speed with \LaTeXe{} for various reasons.

\LaPhYs

\large\LaPhYs

\Large\LaPhYs

\LARGE\LaPhYs

\huge\LaPhYs

\Huge\LaPhYs

%\HUGE\LaPhYs
\normalsize


\mainmatter
% !TEX program = lualatexmk
% !TEX encoding = UTF-8 Unicode

%%%%%%%%%%%%%%%%%%%%%%%%%%%%%%%%%% 80 columns %%%%%%%%%%%%%%%%%%%%%%%%%%%%%%%%%%
% Style Guide:
% DO NOT ALTER THE FIRST TWO LINES OF THIS FILE. Otherwise, lines beginning with
% "% " (note the space after the %) are documentation and may be safely deleted.
% Lines beginning with "%" (no space after the %) are code that may be commented
% or uncommented or deleted to alter functionality.

\chapter{Background Information}

\section{What \TeX{} and \LaTeX{} are Not}
Neither \TeX{} nor \LaTeX{} are word processors. That will be the most
difficult concept for you to grasp. \TeX{} was invented before word
processors became ubiquitous, and you probably did not know that until 
just now. 

\section{What is \TeX?}
\lipsum[5-10]

\section{What is \LaTeX?}
\lipsum[8-12]


% !TEX program = lualatexmk
% !TEX encoding = UTF-8 Unicode

\chapter{Installing a Distribution}

\section{Installing \LaTeX}
Installing a distribution is quite simple, but also dependent upon
your computing platform.

\section{Maintaining Your Distribution}
I highly recommend updating your distribution every day.



\end{document}
